\section{Combining the root orientations}\label{sec:combination}

Proposition~\ref{prop:segment-activity} guarantees activity somewhere on
every long segment across each oriented cut.  It does not yet assert that
the roots are active at the same point.  The wall invariance of the root
indicators supplies the missing relation: on the scale $R_i$, an indicator
can be read as a function of the value of its root alone.  Transitivity
then becomes an approximate addition rule for functions on the real line.
We first obtain activity in every real interval and then use two-edge
cycles to show that the inactive sets have measure zero.

\begin{theorem}[Full activity]\label{thm:full-activity}
For every ordered pair $j\ne k$ in $\{1,\ldots,6\}$,
\begin{equation}\label{eq:full-activity}
 \int_X g_{jk,i}\,\dd\mu_i\longrightarrow1.
\end{equation}
Consequently every root is active almost surely in the limiting state law.
\end{theorem}

\subsection{Functions of a root coordinate}

For each root $\beta$, fix $b_\beta\in L$ such that
$\beta(b_\beta)=1$.  If $x$ is sampled with law $\mu_i$, define the random
measurable functions
\begin{equation}\label{eq:root-line-functions}
 G_{\beta,i}^{x}(s)=g_{\beta,i}\bigl(xa(R_i s b_\beta)\bigr),
 \qquad s\in\R.
\end{equation}
They take values in $\{0,1\}$.  Regard them as elements of the compact
metrizable space
\[
 \mathcal B=\{f\in L^\infty_{\mathrm{loc}}(\R):0\leq f\leq1\},
\]
with weak-* convergence on every bounded interval.  More explicitly,
one may use a countable dense set of $L^1$ test functions on each
$[-q,q]$, $q\in\N$, to metrize this topology.  The finite product
$\mathcal B^{30}$ is also compact and metrizable.  After passing to a
subsequence, the laws of the arrays in \eqref{eq:root-line-functions}
converge to a law of functions $(F_\beta)_\beta$ in this space.
We also write $G_{jk,i}^{x}=G_{\beta_{jk},i}^{x}$ and
$F_{jk}=F_{\beta_{jk}}$.

\begin{lemma}\label{lem:root-line-laws}
For any such convergent subsequence, the limiting law and its prelimit
arrays have the following properties.
\begin{enumerate}
 \item Almost surely, for every $v_0\in L$ with rational coordinates,
 every oriented cut $J$, and every positive rational $\tau$,
 \begin{equation}\label{eq:positive-cut-integral}
  \sum_{\beta\in\mathcal E_J}\int_0^\tau
  F_\beta\bigl(\beta(v_0+t b_J)\bigr)\,\dd t>0.
 \end{equation}
 \item For distinct $j,k,\ell$, and every bounded rectangle
 $B\subset\R^2$, the nonnegative defects satisfy
 \begin{equation}\label{eq:line-transitivity-defect}
  \E_{x\sim\mu_i}\int_B
  G_{jk,i}^{x}(s)G_{k\ell,i}^{x}(t)
  \bigl(1-G_{j\ell,i}^{x}(s+t)\bigr)\,\dd s\,\dd t
  \longrightarrow0.
 \end{equation}
\end{enumerate}
\end{lemma}

\begin{proof}
For $v$ in a fixed bounded subset of $L$, the vector
$v-\beta(v)b_\beta$ belongs to $\ker\beta$ and remains bounded.
Wall invariance \eqref{eq:wall-invariance}, followed by invariance of
$\mu_i$ under the common translation, gives
\begin{equation}\label{eq:common-translate-comparison}
 \sup_{v\text{ in a fixed bounded set}}
 \int_X\left|
 G_{\beta,i}^{x}\bigl(\beta(v)\bigr)
 -g_{\beta,i}\bigl(xa(R_i v)\bigr)
 \right|\,\dd\mu_i(x)\longrightarrow0.
\end{equation}

Fix the parameters in assertion~1.  The left side of
\eqref{eq:positive-cut-integral} is a continuous function of the
weak-* array, because $\beta(b_J)\ne0$ for every crossing root.
By \eqref{eq:common-translate-comparison}, its prelimit differs in
$L^1(\mu_i)$ by $o(1)$ from
\[
 Y_i(x)=\sum_{\beta\in\mathcal E_J}\int_0^\tau
 g_{\beta,i}\bigl(xa(R_i(v_0+t b_J))\bigr)\,\dd t.
\]
Suppose that the limiting observable were zero on a set of probability
$p>0$.  Convergence in law and nonnegativity allow us to choose thresholds
$\varepsilon_i\downarrow0$ and events $E_i$ with
$\mu_i(E_i)\geq p/2$ such that the original prelimit observable is at
most $\varepsilon_i$ on $E_i$.  To see this directly, first choose a
fixed sufficiently small continuity threshold, use convergence in law,
and then decrease the threshold along a diagonal subsequence.
The normalized restrictions $\mu_i|_{E_i}/\mu_i(E_i)$ consequently have
$Y_i$ expectation tending to zero.  Translate these probabilities by
$a(R_i v_0)$ and call the results $\xi_i$.  They satisfy
$\xi_i\leq(2/p)\mu_i$, whereas the change of variable $u=R_i t$
gives
\[
 \frac1{R_i}\int_0^{\tau R_i}\int_X
 A_{J,i}\bigl(xa(u b_J)\bigr)\,\dd\xi_i(x)\,\dd u
 \longrightarrow0.
\]
This contradicts Proposition~\ref{prop:segment-activity}.  The parameters
in assertion~1 range over a countable set, so the assertion holds
simultaneously almost surely.

For assertion~2, the roots $\beta_{jk}$ and $\beta_{k\ell}$ are linearly
independent.  Choose $v(s,t)$ linearly in $(s,t)$ with
\[
 \beta_{jk}(v(s,t))=s,\qquad
 \beta_{k\ell}(v(s,t))=t.
\]
It follows that $\beta_{j\ell}(v(s,t))=s+t$.  On a bounded rectangle,
\eqref{eq:common-translate-comparison} permits replacing all three
functions in the defect by their indicators at the common point
$xa(R_i v(s,t))$, with total expected error $o(1)$.  The expected defect
there is $o(1)$ by the approximate transitivity in
Lemma~\ref{lem:activity-indicators}; its distribution is unchanged by
the translation.  Integrating over the rectangle proves the claim.
\end{proof}

The transitivity defect is nonlinear in the functions, so we do not pass
it directly through weak-* convergence.  Instead we select deterministic
arrays for which the defects themselves tend to zero.

\begin{lemma}\label{lem:deterministic-root-arrays}
If \eqref{eq:full-activity} fails for some root, there is a subsequence
and a deterministic choice of arrays $(G_{\beta,i})_\beta$, each taken
from \eqref{eq:root-line-functions}, with the following properties:
\begin{enumerate}
 \item $G_{\beta,i}\to F_\beta$ weak-* on bounded intervals, and the
 deterministic limiting functions satisfy
 \eqref{eq:positive-cut-integral} for all its stated parameters;
 \item every defect in \eqref{eq:line-transitivity-defect}, without its
 expectation, tends to zero on every bounded rectangle;
 \item for some root $\beta_0$ and some bounded interval $I$ of positive
 length,
 \begin{equation}\label{eq:limiting-deficiency}
  \int_I(1-F_{\beta_0}(s))\,\dd s>0.
 \end{equation}
\end{enumerate}
\end{lemma}

\begin{proof}
Pass to a subsequence on which
$\int_X(1-g_{\beta_0,i})\,\dd\mu_i\geq\varepsilon$ for a fixed
$\varepsilon>0$.  Invariance of $\mu_i$ and Fubini's theorem imply, for
any fixed interval $I$,
\[
 \E\int_I(1-G_{\beta_0,i}^{x}(s))\,\dd s
 =|I|\int_X(1-g_{\beta_0,i})\,\dd\mu_i
 \geq\varepsilon|I|.
\]
The integral is a bounded weak-* continuous observable, so the same
inequality holds for the limiting law.  This law therefore gives positive
probability to a positive deficiency.  Choose a point $F$ in its support
that has a positive deficiency and satisfies assertion~1 of
Lemma~\ref{lem:root-line-laws}.

For each integer $q$, choose a neighborhood $U_q$ of $F$ of diameter
less than $1/q$.  It has positive limiting probability, and thus the
prelimit probabilities of $U_q$ are bounded below by a positive number
for all sufficiently large indices.  There are finitely many root
compositions.  By \eqref{eq:line-transitivity-defect}, the sum of their
defects on $[-q,q]^2$ has expectation tending to zero.  Choose an index
so large that, by Markov's inequality, the probability that this sum
exceeds $1/q$ is less than half the probability of $U_q$.  There is
therefore a sampled array in $U_q$ with total defect at most $1/q$.
Make these choices with increasing indices.  They converge to $F$ and
have vanishing defects on every bounded rectangle, by nonnegativity
and containment in some $[-q,q]^2$.
\end{proof}

\subsection{From cut activity to every interval}

For the rest of the argument, fix deterministic arrays as in
Lemma~\ref{lem:deterministic-root-arrays}.  The labels are still
$\{0,1\}$-valued; their weak-* limits may initially take any value in
$[0,1]$.

\begin{lemma}\label{lem:positive-root-intervals}
For every root $\beta$ and every nonempty open interval $I\subset\R$,
\[
 \int_I F_\beta(s)\,\dd s>0.
\]
\end{lemma}

\begin{proof}
Fix a root $\beta_{jk}$ and a desired interval $I$; replacing it by a
bounded nonempty open subinterval only strengthens what we must prove.
Choose a rational
$v_0\in L$ so that $\beta_{jk}(v_0)$ lies in $I$.  Choose $r>0$ small
enough that
\[
 [\beta_{jk}(v_0)-6r,\ \beta_{jk}(v_0)+6r]\subset I.
\]
For every root $\alpha$, put
$I_\alpha=(\alpha(v_0)-r,\alpha(v_0)+r)$.  Draw the directed edge
$p\to q$ if $\int_{I_{\beta_{pq}}}F_{pq}>0$.

This directed graph is strongly connected.  Indeed, for any nonempty
proper subset $J$, choose a positive rational $\tau$ sufficiently small
that $\alpha(v_0+t b_J)\in I_\alpha$ for every
$\alpha\in\mathcal E_J$ and $0\leq t\leq\tau$.
Equation~\eqref{eq:positive-cut-integral} then supplies an edge from
$J$ to its complement.  A finite directed graph with an outgoing edge
across every nontrivial cut is strongly connected: otherwise the set of
vertices reachable from a suitable vertex is a nonempty proper set
with no outgoing edge.

Take a simple directed path
$j=p_0,p_1,\ldots,p_m=k$, where $m\leq5$.  Write
$\alpha_h=\beta_{p_{h-1}p_h}$ and $I_h=I_{\alpha_h}$.  Weak-*
convergence gives positive lower bounds, for sufficiently large $i$, on
\[
 \int_{I_h}G_{\alpha_h,i}(s)\,\dd s,
 \qquad 1\leq h\leq m.
\]
The product of these integrals is consequently bounded below by a
positive constant.  Interpret it as the measure of tuples
$(s_1,\ldots,s_m)\in\prod_h I_h$ on which all the path labels equal
one.

Except for a set of tuple measure tending to zero, successive composition
of these labels gives
\[
 G_{p_0p_h,i}(s_1+\cdots+s_h)=1,
 \qquad 1\leq h\leq m.
\]
Here is the integration detail.  At step $h$, the failure of this
implication, if all earlier implications held, is bounded by the
transitivity defect for $p_0,p_{h-1},p_h$.  The path is simple, so these
are distinct vertices.  Integration in $s_1,\ldots,s_{h-1}$ pushes
Lebesgue measure on bounded intervals to a measure of bounded density
in their sum.  Multiplying by previous indicator restrictions can only
decrease this measure.  The defect integral on a fixed bounded
rectangle therefore bounds the failure measure by a fixed multiple of
an $o(1)$ quantity.  Summing over the finitely many steps proves the
assertion.  If $m=1$, no composition is necessary.

The final sum $s_1+\cdots+s_m$ lies within $mr$ of
$\alpha_1(v_0)+\cdots+\alpha_m(v_0)=\beta_{jk}(v_0)$, hence in $I$.
Its pushforward density is bounded above by a fixed constant.  Since
a positive measure of tuples has final label one, we obtain
$\liminf_i\int_I G_{jk,i}>0$.  Weak-* convergence proves the lemma.
\end{proof}

We have now located positive activity in every interval for every root.
To eliminate a positive inactive set, we will translate an active set
by sums of parameters from two opposite auxiliary roots.  Composing
successively with those two roots returns to the original root while
adding their two parameters.  Interval positivity supplies parameter
sets of positive measure on both sides of zero.  By mixing normalized
Lebesgue measure on these sets, we will obtain bounded densities of
mean zero, bounded support, and variance bounded away from zero.
The sum of two independent parameters with these densities describes
one cycle increment.  For a fixed bounded interval
$I$, the argument needs a fixed number of independent cycle increments
to have density bounded below on $I-I$: this makes every part of the
active set interact with every part of the inactive set inside $I$.
The next lemma gives that lower bound uniformly over the densities
that arise.

\subsection{A uniform convolution estimate}

\begin{lemma}\label{lem:convolution}
Fix $B,M,v,L>0$.  Consider the probability densities $p$ on $\R$ such
that
\[
 \supp p\subset[-B,B],\qquad 0\leq p\leq M,
 \qquad \int x p(x)\,\dd x=0,
 \qquad \int x^2p(x)\,\dd x\geq v.
\]
There are an integer $m\geq2$ and $c>0$, depending only on
$B,M,v,L$, such that
\[
 p^{*m}(x)\geq c\qquad(|x|\leq L)
\]
for every such density.  The inequality uses the continuous representative
of the convolution density.
\end{lemma}

\begin{proof}
If the family is empty there is nothing to prove.  Otherwise write
$\sigma_p^2=\int x^2p(x)\,\dd x$, so
$v\leq\sigma_p^2\leq B^2$, and let
\[
 \phi_p(t)=\int e^{itx}p(x)\,\dd x.
\]
Bounded support and the zero mean give the uniform expansion
\begin{equation}\label{eq:uniform-characteristic-expansion}
 \phi_p(t)=1-\frac12\sigma_p^2t^2+O(B^3|t|^3).
\end{equation}
Applying the same expansion to
$|\phi_p(t)|^2$, the characteristic function of the difference of two
independent variables with density $p$, shows that for a sufficiently
small $\delta>0$, depending only on $B,v$,
\begin{equation}\label{eq:small-frequency-domination}
 |\phi_p(t)|\leq e^{-vt^2/4}\qquad(|t|\leq\delta).
\end{equation}

There is also a uniform gap below one outside this neighborhood.  For
$|t|\geq\delta$, an angle $\theta$, and $0<a<\pi$, the set
\[
 \{x\in[-B,B]:\operatorname{dist}(tx-\theta,2\pi\Z)<a\}
\]
has Lebesgue measure at most
$a(2B/\pi+4/\delta)$.  Indeed, each complete period contributes the
fraction $a/\pi$ of its length, and at most two incomplete periods
contribute at the ends.  Choose $a$ so small that $M$ times this bound
is at most $1/2$.  With $\theta=\arg\phi_p(t)$, at least half of the
probability lies outside this arc.  Taking the real part after rotation
by $e^{-i\theta}$ gives
\begin{equation}\label{eq:large-frequency-gap}
 \sup_p\sup_{|t|\geq\delta}|\phi_p(t)|
 \leq\frac{1+\cos a}{2}=:\rho<1.
\end{equation}
If $\phi_p(t)=0$, the inequality is immediate and the angle is
irrelevant.  Notice that the estimate holds at arbitrarily large
frequencies; no uniform version of the Riemann--Lebesgue lemma is needed.

Plancherel's identity yields
\[
 \int_\R|\phi_p(t)|^2\,\dd t
 =2\pi\int_\R p(x)^2\,\dd x\leq2\pi M.
\]
For $m\geq2$, $\phi_p^m$ is integrable.  Fourier inversion therefore
defines a continuous representative of $p^{*m}$.  We claim the following
uniform local central limit estimate:
\begin{equation}\label{eq:uniform-local-clt}
 \sup_p\sup_{x\in\R}
 \left|
 \sqrt m\,p^{*m}(x)
 -\frac1{\sqrt{2\pi}\,\sigma_p}
   \exp\left(-\frac{x^2}{2m\sigma_p^2}\right)
 \right|\longrightarrow0.
\end{equation}
To verify it, change variables $u=t\sqrt m$ in the inversion formula.
The difference in \eqref{eq:uniform-local-clt} is at most $1/(2\pi)$
times
\[
 \int_\R\left|
 \phi_p(u/\sqrt m)^m-e^{-\sigma_p^2u^2/2}
 \right|\,\dd u.
\]
On each fixed bounded $u$ interval, the expansion
\eqref{eq:uniform-characteristic-expansion} proves uniform convergence
of the integrand to zero; the logarithm may be taken near one, and its
error is $O(B^3|u|^3/\sqrt m)$.  On
$|u|\leq\delta\sqrt m$, \eqref{eq:small-frequency-domination} bounds
the first term by $e^{-vu^2/4}$, while the Gaussian term is bounded by
$e^{-vu^2/2}$.  Thus the contribution outside a large fixed bounded
interval but inside this region is uniformly small.  On
$|u|>\delta\sqrt m$, \eqref{eq:large-frequency-gap} and Plancherel give
\[
 \int_{|u|>\delta\sqrt m}|\phi_p(u/\sqrt m)|^m\,\dd u
 \leq2\pi M\sqrt m\,\rho^{m-2}\longrightarrow0.
\]
The Gaussian tail there also tends to zero uniformly.  These three
estimates prove \eqref{eq:uniform-local-clt}.

For $|x|\leq L$ and all sufficiently large $m$, the Gaussian term in
\eqref{eq:uniform-local-clt} is at least
$1/(2\sqrt{2\pi}B)$, uniformly in $p$.  Increase $m$ so that the
error is at most $1/(4\sqrt{2\pi}B)$.  Then
\[
 p^{*m}(x)\geq\frac1{4\sqrt{2\pi}B\sqrt m}
 \quad(|x|\leq L),
\]
which is the required conclusion.
\end{proof}

\subsection{Cycles remove the inactive sets}

\begin{proof}[Proof of Theorem~\ref{thm:full-activity}]
Suppose the theorem fails, and take the deterministic arrays and limiting
functions from Lemma~\ref{lem:deterministic-root-arrays}.  Fix an edge
$j\to k$, and choose $\ell$ distinct from $j,k$.
Lemma~\ref{lem:positive-root-intervals} supplies, for each of the roots
$\ell j$ and $j\ell$, positive limiting integral on each of the
intervals $(-2,-1)$ and $(1,2)$.  Accordingly, for all sufficiently
large $i$, each of the four sets where the relevant $G$ equals one
has length bounded below by a positive constant.

For each auxiliary root, normalize Lebesgue measure on its two sets
to probability densities $p_i^-$ and $p_i^+$.  Their means belong to
$[-2,-1]$ and $[1,2]$, respectively.  There is a convex combination
\[
 p_i=\lambda_i p_i^-+(1-\lambda_i)p_i^+
\]
with mean zero.  The coefficients can be chosen in $[1/3,2/3]$ by
solving the mean-zero equation.  Thus these densities have support in
$[-2,2]$, a common finite upper bound, and variance at least one.
They are supported entirely where the corresponding auxiliary
indicator is one.  Denote the density for $\ell j$ by $p_i$ and that
for $j\ell$ by $q_i$.

Starting from $G_{jk,i}(t)=1$, composition with a parameter $s$ in the
support of $p_i$ gives the edge $\ell k$, and composition with a
parameter $s'$ in the support of $q_i$ returns to $jk$:
\[
 \ell j+jk=\ell k,\qquad j\ell+\ell k=jk.
\]
Both are root compositions with three distinct vertices.  Their defect
estimates imply, for every fixed bounded interval $I$,
\begin{equation}\label{eq:one-cycle-error}
 \int_I\int_\R\int_\R
 G_{jk,i}(t)\bigl(1-G_{jk,i}(t+s+s')\bigr)
 p_i(s)q_i(s')\,\dd s\,\dd s'\,\dd t\longrightarrow0.
\end{equation}
Indeed, except on null sets of parameters, both auxiliary labels are
one.  The indicator in \eqref{eq:one-cycle-error} is then bounded by
the sum of the two composition defects.  For the first defect,
integrate out $s'$ and use the uniform density bound on $p_i$.
For the second, the variable $t+s$ has a bounded density on a fixed
bounded interval when $t\in I$ is integrated with Lebesgue measure
and $s$ with $p_i$.  The bound on $q_i$ and the defect estimate for
$(j\ell,\ell k)$ finish the verification.

Put $r_i=p_i*q_i$.  These cycle-increment densities have support in
$[-4,4]$, a common density bound, mean zero, and variance at least two.
We can iterate \eqref{eq:one-cycle-error} any fixed number $m$ of times
to get
\begin{equation}\label{eq:repeated-cycle-error}
 \int_I\int_\R
 G_{jk,i}(t)\bigl(1-G_{jk,i}(t+u)\bigr)
 r_i^{*m}(u)\,\dd u\,\dd t\longrightarrow0.
\end{equation}
For clarity, along a sequence of $m$ increments, a path starting with
label one and ending with label zero has a first transition from one
to zero.  Bound its indicator by the sum of the indicators of these
transitions.  The intermediate positions lie in
$I+[-4m,4m]$.  After integrating all preceding increments, their
distributions have densities bounded by one relative to Lebesgue
measure on that larger interval, because the initial $t$ is integrated
over $I$ and all increments have probability laws.  Each term is
therefore bounded by the one-cycle error on the larger interval.
There are only $m$ terms, proving \eqref{eq:repeated-cycle-error}.

Fix a bounded interval $I$ of positive length and apply
Lemma~\ref{lem:convolution} to the $r_i$, with $L$ large enough that
$I-I\subset[-L,L]$.  Choose its fixed power $m$ and lower bound
$c>0$.  In \eqref{eq:repeated-cycle-error}, retain only the pairs
$t,t+u\in I$.  We obtain
\[
 c\left(\int_I G_{jk,i}(t)\,\dd t\right)
   \left(\int_I(1-G_{jk,i}(t))\,\dd t\right)
 \longrightarrow0.
\]
The first factor has positive lower limit by
Lemma~\ref{lem:positive-root-intervals}.  Hence
$\int_I(1-G_{jk,i})\to0$ for every such interval.  Weak-* convergence
shows that $F_{jk}=1$ almost everywhere.  The argument applies to every
root, contradicting the deficiency
\eqref{eq:limiting-deficiency}.  Thus \eqref{eq:full-activity} holds.

By Lemma~\ref{lem:activity-indicators}, the sequence indicators converge
jointly in law to the activity indicators of the limiting state.  Since
each has expectation tending to one, every limiting root indicator
equals one almost surely.  There are only thirty roots, so this holds
for all of them simultaneously.
\end{proof}
